% بخش ۸: توابع هزینه و معیارهای بهینه‌سازی
\section{توابع هزینه و معیارهای بهینه‌سازی}
\label{sec:loss_functions}

طراحی تابع هزینه در مسئلهٔ تخمین عمق تک‌دوربینه، نقش تعیین‌کننده‌ای در نوع اطلاعاتی دارد که مدل در فرآیند آموزش فرا می‌گیرد. از آنجا که هدف این پژوهش، دستیابی به تخمین عمق \emph{متریک پایدار} و \emph{ناوابسته به تغییر پارامترهای درونی دوربین} است، انتخاب تابع هزینه به‌گونه‌ای انجام شده که مدل بتواند مقیاس مطلق عمق را یاد بگیرد و در عین حال نسبت به تغییرات پارامترهای دوربین پایدار باشد.

در این پژوهش، بسته به نوع مدل (متریک یا نسبی)، از توابع هزینهٔ مناسبی استفاده شده است که در ادامه تشریح می‌شوند.

\subsection{هزینهٔ SiLog برای مدل‌های متریک}

برای آموزش مدل متریک که هدف آن تولید تخمین عمق مطلق بر حسب متر است، از تابع هزینهٔ لگاریتمی تغییرناپذیر نسبت به مقیاس\footnote{\lr{Scale-Invariant Logarithmic Loss}} استفاده می‌شود. این تابع هزینه که به اختصار SiLog نامیده می‌شود، به‌گونه‌ای طراحی شده که به‌جای تفاوت مستقیم مقادیر عمق، بر تفاوت لگاریتمی آن‌ها تمرکز دارد:

\begin{equation}
\mathcal{L}_{\text{SiLog}} = \sqrt{\frac{1}{N} \sum_{i=1}^{N} d_i^2 - \frac{\lambda}{N^2} \left( \sum_{i=1}^{N} d_i \right)^2},
\end{equation}

که در آن $N$ تعداد پیکسل‌های دارای برچسب معتبر، $d_i = \log D_i^{GT} - \log D_i^{pred}$ تفاوت لگاریتمی بین مقدار واقعی و مقدار پیش‌بینی‌شده در پیکسل $i$-ام است. ضریب $\lambda$ معمولاً برابر $0.5$ در نظر گرفته می‌شود تا تعادلی بین خطای مطلق و خطای نسبی برقرار گردد.

ویژگی مهم SiLog آن است که ضمن حفظ حساسیت به مقیاس مطلق عمق، باعث می‌شود خطا در فواصل دور و نزدیک به‌صورت متعادل‌تری در نظر گرفته شود. استفاده از فضای لگاریتمی، مدل را نسبت به خطاهای بزرگ در عمق‌های دور مقاوم‌تر می‌کند و از تمرکز بیش از حد بر خطاهای نقطه‌ای بسیار بزرگ جلوگیری می‌نماید.

\subsection{هزینهٔ تغییرناپذیر نسبت به مقیاس و جابجایی برای مدل‌های نسبی}

برای آموزش مدل در حالت نسبی (basic mode) که خروجی آن دارای ابهام مقیاس است، از تابع هزینهٔ تغییرناپذیر نسبت به مقیاس و جابجایی\footnote{\lr{Scale-Shift Invariant Loss}} استفاده می‌شود. این تابع هزینه ابتدا پیش‌بینی مدل را با استفاده از روش کمترین مربعات بر مقدار واقعی منطبق می‌کند:

\begin{equation}
D^*_{pred} = s \cdot D_{pred} + t,
\end{equation}

که در آن $s$ و $t$ مقیاس و جابجایی بهینه هستند که به‌صورت تحلیلی از طریق کمینه‌سازی خطای پیکسل‌به-پیکسل محاسبه می‌شوند.

تابع هزینهٔ نهایی ترکیبی از خطای مطلق پیکسل‌ها ($L_1$) و خطای تطابق گرادیان است:

\begin{equation}
\mathcal{L}_{\text{SSI}} = (1 - \alpha) \mathcal{L}_{\text{pixel}} + \alpha \mathcal{L}_{\text{grad}},
\end{equation}

که در آن $\mathcal{L}_{\text{pixel}}$ خطای $L_1$ میان تخمین منطبق‌شده و عمق واقعی است و $\mathcal{L}_{\text{grad}}$ خطای تطابق گرادیان در چند مقیاس را اندازه می‌گیرد. پارامتر $\alpha$ (معمولاً 0.5) وزن نسبی این دو مؤلفه را تعیین می‌کند. هزینهٔ گرادیان باعث می‌شود که لبه‌های عمق و ناپیوستگی‌ها در تصویر با دقت بیشتری بازسازی شوند.

\subsection{تمایز میان توابع هزینهٔ آموزش و معیارهای ارزیابی}

نکته‌ای حائز اهمیت در این پژوهش، تمایز آگاهانه میان توابع هزینهٔ مورد استفاده در آموزش و معیارهای به‌کاررفته در ارزیابی مدل است. این تمایز از آن جهت ضروری است که اهداف آموزش و ارزیابی لزوماً یکسان نیستند.

\subsubsection{توابع هزینه در آموزش}

در مرحلهٔ آموزش:
\begin{itemize}
    \item برای مدل‌های متریک از \textbf{SiLog} استفاده می‌شود
    \item برای مدل‌های نسبی از \textbf{Scale-Shift Invariant Loss} استفاده می‌شود
\end{itemize}

\subsubsection{معیارهای ارزیابی}

در مرحلهٔ ارزیابی، از معیارهای متنوع‌تری بسته به نوع مدل استفاده می‌شود:

برای مدل‌های متریک، از معیارهای استاندارد تخمین عمق استفاده می‌گردد:
\begin{itemize}
    \item \textbf{خطای مطلق نسبی (AbsRel):} 
    \begin{equation}
    \text{AbsRel} = \frac{1}{N}\sum_{i=1}^{N} \frac{|D_i^{pred} - D_i^{GT}|}{D_i^{GT}}
    \end{equation}
    
    \item \textbf{RMSE:}
    \begin{equation}
    \text{RMSE} = \sqrt{\frac{1}{N}\sum_{i=1}^{N} (D_i^{pred} - D_i^{GT})^2}
    \end{equation}
    
    \item \textbf{معیارهای دقت آستانه‌ای ($\delta_i$):} درصد پیکسل‌هایی که $\max\left(\frac{D^{pred}}{D^{GT}}, \frac{D^{GT}}{D^{pred}}\right) < 1.25^i$ برای $i \in \{1, 2, 3\}$
\end{itemize}

برای مدل‌های نسبی که دارای ابهام مقیاس هستند، ابتدا یک تراز مقیاس (معمولاً از طریق میانه) اعمال شده و سپس معیارهای فوق محاسبه می‌شوند. همچنین معیار SiLog که ذاتاً تغییرناپذیر نسبت به مقیاس است، برای مقایسهٔ این نوع مدل‌ها مناسب‌تر است.

\subsection{نقش پارامترهای درونی دوربین در آموزش}

نقش ماتریس مشخصهٔ دوربین در این چارچوب، نه از طریق افزودن یک مؤلفهٔ مستقیم به تابع هزینه، بلکه از طریق \emph{تزریق اطلاعات هندسی به معماری مدل} ایفا می‌شود. با در اختیار داشتن اطلاعات صریح پارامترهای درونی، مدل قادر است تخمین‌های متریک خود را با هندسهٔ واقعی تصویربرداری تطبیق دهد.

در مدل پیشنهادی، پارامترهای درونی دوربین از طریق یک واحد کدگذار مجزا پردازش شده و به‌صورت توکن دوربین به لایه‌های میانی Vision Transformer تزریق می‌شوند. این اطلاعات به مدل امکان می‌دهد تا در فرآیند یادگیری، رابطهٔ میان ویژگی‌های تصویری و مقیاس واقعی عمق را با در نظر گرفتن مشخصات هندسی دوربین یاد بگیرد.

به این ترتیب، تابع هزینهٔ SiLog که ذاتاً به مقیاس مطلق عمق حساس است، می‌تواند به‌درستی مدل را برای تولید تخمین‌های متریک پایدار هدایت کند. این پایداری در مواجهه با داده‌هایی که تنها در پارامترهای درونی دوربین تفاوت دارند، در آزمایش‌های تجربی به‌وضوح مشاهده می‌شود.

\subsection{حساسیت RMSE به پارامترهای دوربین و کاربرد آن در ارزیابی}

معیار RMSE که در این پژوهش برای ارزیابی مدل‌های متریک استفاده شده، به‌صورت مستقیم به مقیاس عمق وابسته است:

\begin{equation}
\text{RMSE} = \sqrt{\frac{1}{N}\sum_{i=1}^{N} (D_i^{pred} - D_i^{GT})^2}.
\end{equation}

این وابستگی باعث می‌شود RMSE نسبت به تغییر پارامترهای درونی دوربین، به‌ویژه فاصلهٔ کانونی، حساس باشد. این ویژگی، هرچند در نگاه اول یک محدودیت به نظر می‌رسد، اما در واقع ابزار مناسبی برای آشکارسازی ناپایداری متریک مدل‌های عمق محسوب می‌شود.

در این پژوهش، از این حساسیت به‌عنوان معیاری برای سنجش پایداری مدل استفاده شده است. مدلی که نسبت به تغییر پارامترهای دوربین پایدار باشد، باید تغییرات کمی در RMSE از خود نشان دهد، حتی زمانی که تصاویر با مشخصات هندسی متفاوت تصویربرداری شده باشند.

\subsection{جمع‌بندی}

در این بخش، توابع هزینهٔ مورد استفاده در آموزش مدل‌های متریک و نسبی تشریح شد. استفاده از SiLog برای مدل‌های متریک، ضمن حفظ حساسیت به مقیاس مطلق عمق، مدل را نسبت به خطاهای بزرگ مقاوم می‌کند. تزریق پارامترهای درونی دوربین به معماری مدل، امکان یادگیری رابطهٔ صحیح میان ویژگی‌های تصویری و عمق واقعی را فراهم می‌آورد. در فصل بعد، نتایج تجربی حاصل از این طراحی و تحلیل آماری پایداری مدل نسبت به تغییر مشخصات دوربین ارائه خواهد شد.


% بخش ۹: جزئیات پیاده‌سازی
\section{جزئیات پیاده‌سازی}
\label{sec:implementation_details}

در این بخش، جزئیات فنی مربوط به پیاده‌سازی مدل پیشنهادی، تنظیمات آموزش، پیش‌پردازش داده‌ها، ساختار کد و ملاحظات عملی لازم برای بازتولید نتایج ارائه می‌شود. هدف از این بخش، مستندسازی دقیق تمامی تصمیمات پیاده‌سازی است تا امکان تکرار آزمایش‌ها و تحلیل نتایج برای پژوهشگران دیگر فراهم شود.

\subsection{چارچوب نرم‌افزاری و زیرساخت محاسباتی}

تمامی پیاده‌سازی‌ها در چارچوب \lr{PyTorch} (نسخهٔ 2.0 و بالاتر) انجام شده و کد پایه بر اساس مخزن رسمی \lr{Depth Anything V2} توسعه داده شده است. معماری کلی شامل یک انکودر مبتنی بر \lr{DINOv2 Vision Transformer} و یک دیکودر از خانوادهٔ \lr{Dense Prediction Transformer (DPT)} است.

آموزش مدل‌ها بر روی پردازنده‌های گرافیکی \lr{NVIDIA} با معماری Ampere یا جدیدتر انجام شده و تمامی آزمایش‌ها به‌صورت تک‌کارت گرافیک و با دقت ممیز شناور ۳۲ بیتی اجرا شده‌اند. برای تضمین تکرارپذیری، بذر تصادفی\footnote{\lr{Random Seed}} برای کتابخانه‌های PyTorch، NumPy و Python در ابتدای هر آزمایش ثابت (معمولاً بر روی مقدار 42) تنظیم شده است.

\subsection{پیش‌پردازش داده‌ها}

تمامی تصاویر ورودی پیش از ورود به شبکه، به‌صورت یکنواخت پیش‌پردازش می‌شوند. تصاویر ابتدا به اندازهٔ ثابت $518 \times 518$ پیکسل تغییر اندازه داده شده و سپس در صورت نیاز، برش‌های تصادفی با اندازهٔ کوچک‌تر برای افزایش تنوع داده‌ای در مرحلهٔ آموزش اعمال می‌گردد. این راهبرد، مطابق با تنظیمات مدل پایه \lr{Depth Anything V2}، به بهبود تعمیم‌پذیری مدل کمک می‌کند \cite{DepthAnythingV2}.

مقادیر پیکسل تصاویر به بازهٔ $[0,1]$ نرمال‌سازی شده و سپس با استفاده از میانگین و واریانس داده‌های \lr{ImageNet} استانداردسازی می‌شوند:

\begin{equation}
I_{norm} = \frac{I - \mu_{ImageNet}}{\sigma_{ImageNet}},
\end{equation}

که در آن $\mu_{ImageNet} = [0.485, 0.456, 0.406]$ و $\sigma_{ImageNet} = [0.229, 0.224, 0.225]$ هستند. این مرحله به‌ویژه برای سازگاری با وزن‌های ازپیش‌آموزش‌دیدهٔ انکودر \lr{DINOv2} ضروری است.

\subsection{مدیریت و نرمال‌سازی پارامترهای درونی دوربین}

برای هر تصویر ورودی، ماتریس مشخصهٔ دوربین $K$ استخراج شده و پارامترهای اصلی آن شامل فاصلهٔ کانونی در راستاهای افقی و عمودی ($f_x, f_y$) و مختصات نقطهٔ اصلی ($c_x, c_y$) مورد استفاده قرار می‌گیرند.

به‌منظور جلوگیری از ناپایداری عددی و تسهیل فرآیند یادگیری، این پارامترها پیش از ورود به شبکه نرمال‌سازی می‌شوند:

\begin{equation}
\tilde{f}_x = \frac{f_x}{W}, \quad
\tilde{f}_y = \frac{f_y}{H}, \quad
\tilde{c}_x = \frac{c_x}{W}, \quad
\tilde{c}_y = \frac{c_y}{H},
\end{equation}

که در آن $W$ و $H$ به‌ترتیب عرض و ارتفاع تصویر ورودی هستند.

بردار نهایی پارامترهای دوربین به‌صورت

\begin{equation}
\mathbf{k} = [\tilde{f}_x, \tilde{f}_y, \tilde{c}_x, \tilde{c}_y]
\end{equation}

تعریف شده و سپس از طریق یک شبکهٔ تمام‌متصل کوچک (شامل دو لایهٔ خطی با تابع فعال‌سازی ReLU) به یک نمایش نهان $\mathbf{z}_k$ با بُعد 256 نگاشت می‌شود. این نمایش نهان در مرحلهٔ دیکودینگ به ویژگی‌های تصویری تزریق می‌گردد.

\subsection{تنظیمات آموزش و راهبرد بهینه‌سازی}

مدل پیشنهادی با استفاده از بهینه‌ساز \lr{AdamW} \cite{Loshchilov2019AdamW} آموزش داده شده است. نرخ یادگیری پایه برای انکودر \lr{DINOv2} برابر با $5 \times 10^{-6}$ در نظر گرفته شده، در حالی که برای دیکودر و لایه‌های افزوده‌شده (شامل کدگذار دوربین و سر تخمین عمق)، از یک ضریب تقویتی (معمولاً 10 برابر) استفاده شده که نرخ یادگیری این بخش‌ها را به $5 \times 10^{-5}$ می‌رساند.

پارامترهای بهینه‌ساز به‌صورت زیر تنظیم شده‌اند:

\begin{itemize}
    \item ضرایب ممنتوم: $\beta_1 = 0.9$، $\beta_2 = 0.999$
    \item ضریب کاهش وزن (\lr{Weight Decay}): $0.01$
    \item پارامتر $\epsilon$ برای پایداری عددی: $10^{-8}$
\end{itemize}

در راستای جلوگیری از بیش‌برازش و حفظ دانش ازپیش‌آموخته‌شده، وزن‌های انکودر \lr{DINOv2} در طول آموزش فریز شده‌اند. این تنظیم، مطابق با گزینهٔ پیش‌فرض \lr{\texttt{--freeze-dinov2}} در اسکریپت آموزش اعمال شده است. فریز کردن انکودر باعث می‌شود تنها لایه‌های سر شبکه\footnote{\lr{Head Layers}} و واحد کدگذار دوربین آموزش ببینند، که این امر از واگرایی نمایش‌های ویژگی جلوگیری کرده و همگرایی سریع‌تر را تضمین می‌کند.

برای زمان‌بندی نرخ یادگیری، از یک استراتژی ترکیبی شامل \lr{warmup} خطی در ابتدای آموزش و سپس کاهش چندجمله‌ای (\lr{Polynomial Decay}) با توان $0.9$ استفاده شده است. تعداد دوره‌های آموزش معمولاً بین 20 تا 40 epoch بوده و اندازهٔ دسته بسته به حافظهٔ موجود بین 2 تا 8 تنظیم شده است. در صورت محدودیت حافظه، از تجمع گرادیان\footnote{\lr{Gradient Accumulation}} برای افزایش اندازهٔ دستهٔ مؤثر استفاده شده است.

\subsection{دیتاست‌های مورد استفاده در پیاده‌سازی}

پیاده‌سازی مدل بر روی مجموعه‌ای متنوع از دیتاست‌های داخلی و خارجی انجام شده است که شامل صحنه‌های indoor و outdoor با مشخصات تصویربرداری متفاوت هستند. از جمله دیتاست‌های مورد استفاده می‌توان به موارد زیر اشاره کرد:

\begin{itemize}
    \item \textbf{\lr{VKITTI2}:} دیتاست شبیه‌سازی‌شده برای صحنه‌های رانندگی با ground truth دقیق عمق و پارامترهای دوربین
    \item \textbf{\lr{KITTI}:} دیتاست واقعی رانندگی خودکار با اطلاعات عمق از لایدار
    \item \textbf{\lr{NYUv2}:} دیتاست صحنه‌های داخلی با عمق حاصل از سنسور Kinect
    \item \textbf{\lr{Hypersim}:} دیتاست سنتتیک با تنوع بالا در صحنه‌های داخلی و کنترل کامل بر پارامترهای دوربین
\end{itemize}

این تنوع دیتاست‌ها، به‌ویژه از منظر تفاوت در پارامترهای درونی دوربین و میدان دید، نقش مهمی در ارزیابی پایداری متریک مدل پیشنهادی ایفا می‌کند. برای هر دیتاست، فایل‌های لیست تقسیم‌بندی داده (train/validation/test splits) تهیه شده که شامل مسیر تصاویر، نقشه‌های عمق و پارامترهای درونی دوربین در قالب فایل‌های \lr{.npy} هستند.

\subsection{ساختار کد و اسکریپت‌های اصلی}

ساختار کد به‌گونه‌ای طراحی شده که اجزای اصلی شامل ماژول‌های بارگذاری داده، تعریف مدل، آموزش و ارزیابی به‌صورت مجزا پیاده‌سازی شوند. ساختار کلی پروژه به شرح زیر است:

\begin{verbatim}
project_root/
├── models/
│   ├── raw_models/
│   │   ├── DepthAnythingV2/         # مدل پایه
│   │   ├── DepthAnythingV2-revised/ # مدل با intrinsics
│   │   └── Depth-Anything-3/
│   └── base.py                      # کلاس‌های پایه
├── datasets/
│   ├── base.py                      # کلاس پایه دیتاست
│   ├── training_datasets.py        # دیتاست‌های آموزش
│   └── (middlebury, cityscapes, ...)
├── src/
│   ├── pipeline.py                  # خط لوله اصلی
│   ├── metrics.py                   # معیارهای ارزیابی
│   └── visualization.py             # ابزارهای تجسم
├── train.py                         # اسکریپت اصلی آموزش
├── evaluate.py                      # اسکریپت ارزیابی
└── compare_models.py                # مقایسه آماری
\end{verbatim}

اسکریپت \lr{\texttt{train.py}} مسئول اجرای فرآیند آموزش بوده و از طریق آرگومان‌های ورودی، امکان تنظیم پارامترهای زیر را فراهم می‌کند:

\begin{itemize}
    \item \texttt{--use-camera-intrinsics}: فعال‌سازی استفاده از پارامترهای دوربین
    \item \texttt{--model-type}: انتخاب نوع مدل (metric یا basic)
    \item \texttt{--encoder}: انتخاب اندازهٔ انکودر (vits, vitb, vitl, vitg)
    \item \texttt{--use-distillation}: فعال‌سازی تقطیر دانش
    \item \texttt{--freeze-dinov2}: فریز کردن وزن‌های انکودر
    \item \texttt{--lambda-distill} و \texttt{--lambda-metric}: تنظیم ضرایب تابع هزینه
\end{itemize}

علاوه بر اسکریپت آموزش، اسکریپت‌های مستقلی برای مقایسهٔ خروجی مدل‌ها، تحلیل آماری نتایج (شامل آزمون t زوجی و محاسبهٔ فواصل اطمینان bootstrap) و تجسم نتایج پیاده‌سازی شده‌اند.

\subsection{ملاحظات عملی و محدودیت‌ها}

در پیاده‌سازی عملی، برخی محدودیت‌ها و چالش‌ها شناسایی و مدیریت شده‌اند:

\begin{itemize}
    \item \textbf{محدودیت حافظهٔ GPU:} با توجه به حجم بالای مدل‌های Vision Transformer، از تکنیک‌هایی نظیر gradient accumulation و mixed precision training استفاده شده است.
    
    \item \textbf{عدم دسترسی به پارامترهای دقیق دوربین:} برای برخی دیتاست‌ها، پارامترهای دوربین از روی اطلاعات EXIF تصاویر یا مستندات دیتاست استخراج یا تخمین زده شده است. در مواردی که اطلاعات دقیق در دسترس نبوده، از مقادیر نموذج (مانند $f = W$ برای تصاویر با FoV استاندارد) استفاده شده است.
    
    \item \textbf{تنوع در کیفیت ground truth:} دیتاست‌های مختلف دارای سطوح متفاوتی از دقت در نقشه‌های عمق مرجع هستند. برای رفع این مسئله، در فرآیند ارزیابی از masking برای حذف پیکسل‌های نامعتبر استفاده شده است.
    
    \item \textbf{تکرارپذیری:} علی‌رغم تلاش برای ثابت نگه‌داشتن تمامی بذرهای تصادفی، برخی عملیات GPU ممکن است غیرقطعی باشند. برای کاهش این اثر، هر آزمایش حداقل سه بار تکرار شده و میانگین نتایج گزارش شده است.
\end{itemize}

\subsection{جمع‌بندی}

در این بخش، جزئیات پیاده‌سازی مدل پیشنهادی شامل چارچوب نرم‌افزاری، پیش‌پردازش داده‌ها، مدیریت پارامترهای درونی دوربین، تنظیمات آموزش، دیتاست‌های مورد استفاده و ساختار کد تشریح شد. این توضیحات، مبنای فنی لازم برای درک و تحلیل نتایج تجربی ارائه‌شده در فصل بعد را فراهم می‌کنند و امکان بازتولید آزمایش‌ها توسط سایر پژوهشگران را تسهیل می‌نمایند.
